% ==============================================================================
%  FRONT MATTER: declaration, abstract, acknowledgements
% ==============================================================================

\section*{Declaration}
\addcontentsline{toc}{section}{Declaration}

I declare that this report is my own work and that, except where reference is made to
the work of others, it has not been submitted in whole or in part for any other degree
or qualification at this or any other university.

\vspace{1.5cm}
\begin{tabular}{@{}l@{\hspace{3cm}}l@{}}
\rule{5cm}{0.4pt} & \rule{4cm}{0.4pt} \\
Akash Abeywickrame (S/20/302) & Date \\[1.5cm]
\rule{5cm}{0.4pt} & \rule{4cm}{0.4pt} \\
Dr. Hakim Usoof (Supervisor) & Date \\
\end{tabular}

\clearpage
\section*{Abstract}
\addcontentsline{toc}{section}{Abstract}

Camera glasses record everyone in front of the wearer, including strangers who never
agreed to be filmed. The two protections in use today both fall short: a small
recording light that bystanders rarely notice, and blurring every face, which also
erases the friends and companions the wearer is actually with. The one existing system
that blurs selectively depends on eye tracking and a microphone array that ordinary
camera glasses do not have.

In this project I designed and evaluated a pipeline that blurs bystanders but leaves the
wearer's interaction partners visible, using only the glasses' RGB camera and without
recognising anyone. Faces are detected with YOLOv8n-face, followed over time with a
motion-only SORT tracker, and described by four cues: how much of the frame a face
covers, how close it is to the centre of view, which way the head is turned, and
whether the mouth is moving. A rule-based classifier with temporal confirmation and
latching turns these cues into a partner or bystander decision for each face. For
recorded video, an offline second pass decides once per person for the whole video: it
keeps a confirmed partner visible for their whole time on screen, keeps people who stay with
the wearer's group, and rejoins a person who briefly leaves the view, all without
recognising anyone.

I evaluated the system on 21 first-person videos, with interaction scenes from the
EgoBlur release and bystander scenes from YouTube (131,187 frames, 2,736 face tracks),
against 197 labelled people from three human labelling studies. In the
largest, five annotators labelled 102 short video clips with almost perfect agreement
(Fleiss' $\kappa = 0.887$). The live system blurred 96.97\% of bystanders and kept
56.52\% of partners visible; when it kept a face visible, it was right 97.5\% of the
time. The offline second pass raised partner preservation to 85.51\% while still
blurring 90.91\% of bystanders. An ablation showed that face size drives most
decisions, and letting the head-pose cue lower the size threshold for faces turned towards
the wearer raised live partner preservation to 78.26\%, with 90.91\% of bystanders still
blurred, when tested on videos it was not tuned on. The classifier
itself takes under a millisecond per frame. With the face detector exported to OpenVINO,
the whole pipeline runs in real time on a laptop CPU, at 31 frames per second without and
22 with the landmark cues, while still finding 99.7\% of the faces. Projections for the
Ray-Ban Meta glasses reach about 18 frames per second when detection runs on their neural
processing unit.

\textbf{Keywords:} bystander privacy, augmented reality, egocentric vision, selective
anonymization, face detection, edge computing.

\clearpage
\section*{Acknowledgements}
\addcontentsline{toc}{section}{Acknowledgements}

I am grateful to my supervisor, Dr. Hakim Usoof, for his guidance and for detailed
feedback that shaped this report. I thank the Department of Statistics and Computer
Science, University of Peradeniya, and the friends who gave their time to label the
evaluation clips. % TODO(Akash): add anyone else you wish to thank.
